Streaming the layers of a model through a small GPU-memory ring is a max-plus linear system, and its iteration time is the maximum mean of a family of tilings of the layer sequence. For layers of equal cost this gives a complete answer to the placement question. When a layer computes at least as long as its copy takes, two slots suffice and resident layers are useless. Otherwise, spreading the resident layers evenly attains the lower bound $\max(Lc,mf)$ as soon as $(k-1)f\ge2c$, while a contiguous block, the result of filling GPU memory in model order, takes $\max(Lc,mf,rc+2c+(m-k)f)$, up to twice as long. On a Tesla T4 with page-locked host memory, the same bandwidth class as current CXL expanders, the model predicted \nAccN{} measured iteration times with a mean error of \nAccMape\%, and even spreading was up to \nMeasPenaltyMax$\times$ faster than the contiguous prefix. The rules are cheap enough for any offloading runtime to adopt; the open questions are tensor-level units, streamed key--value caches and power-limited operation, in which the per-layer times themselves depend on the schedule.
